\chapter{O padrão-ouro e os antecedentes da Caixa de Conversão}

\section{O padrão-ouro como instituição internacional}

Para compreender o debate brasileiro de 1906, é necessário separar o padrão-ouro como conjunto de regras, como prática histórica e como linguagem de legitimidade. Em sua formulação mais simples, o regime definia a unidade monetária por um peso de ouro, assegurava a conversão entre moeda e metal à paridade estabelecida e permitia o movimento internacional do ouro. A adoção simultânea dessas regras por diferentes países produzia taxas de câmbio fixas entre suas moedas, sujeitas aos custos de transporte, seguro e operação do metal.

Oliveira e Silva sintetizam quatro requisitos: definição do preço da moeda em ouro, garantia de conversibilidade, manutenção de reservas compatíveis com as necessidades externas e ausência de restrições à entrada e à saída do metal \cite{oliveirasilva2001}. Nessa descrição, um superávit no balanço de pagamentos acrescentaria ouro às reservas, possibilitaria expansão monetária e tenderia a reduzir juros ou elevar preços. Um déficit provocaria saída de reservas, contração da moeda, aumento dos juros ou queda de preços. As mudanças resultantes na demanda, nos preços relativos e nos fluxos de capital deveriam corrigir o desequilíbrio externo.

Esse mecanismo é uma representação analítica, não a descrição completa do que bancos centrais e governos faziam. Autoridades podiam alterar taxas de redesconto, manejar ativos externos, obter empréstimos e cooperar para adiar ou distribuir o ajuste. A própria credibilidade da paridade afetava o sentido dos fluxos. Se investidores acreditassem que uma taxa seria preservada, a elevação dos juros poderia atrair capitais antes que uma grande quantidade de ouro deixasse o país. Se esperassem desvalorização ou suspensão, a retirada preventiva de recursos tornava a defesa do regime mais difícil.

A literatura que interpreta o padrão-ouro como regra enfatiza essa dimensão expectacional. O compromisso com a paridade podia restringir decisões discricionárias e sinalizar a disposição de preservar contratos denominados em ouro \cite{bordokydland1990}. Isso não significa que a regra fosse absoluta. Bordo e Schwartz propõem entendê-la como contingente: a conversibilidade deveria ser mantida em circunstâncias normais, mas suspensões diante de choques extraordinários podiam ser compatíveis com a reputação do regime se fossem acompanhadas da expectativa de retorno posterior \cite{bordoschwartz1994}. A história do padrão-ouro, portanto, inclui tanto estabilidade quanto crises, empréstimos de socorro, controles temporários e suspensões.

Eichengreen chama atenção para o suporte político e internacional necessário ao funcionamento do sistema \cite{eichengreen1996}. A prioridade conferida à paridade não era apenas uma técnica monetária. Dependia de coalizões dispostas a aceitar seus custos internos e da capacidade de cooperação entre autoridades. À medida que a participação política e a pressão por emprego e renda se ampliaram, tornou-se mais difícil subordinar a política doméstica ao ajuste externo sem contestação. Mesmo antes de 1914, essa tensão se manifestava de forma desigual entre países.

\subsection{Centro, periferia e assimetria do ajuste}

O modelo automático também obscurece a posição diferenciada das economias centrais e periféricas. Países exportadores de poucos produtos primários estavam expostos a oscilações de preço, colheita e demanda que podiam alterar rapidamente suas receitas externas. Seus mercados financeiros internos eram menos profundos, e a oferta de capitais internacionais podia desaparecer justamente quando a queda das exportações aumentava a necessidade de financiamento. O ajuste por saída de ouro e contração monetária tendia, nessas condições, a reforçar a desaceleração já iniciada pelo setor externo.

Franco analisa essas assimetrias sistêmicas e questiona a ideia de que o mecanismo distribuía de modo equivalente os encargos entre países superavitários e deficitários \cite{franco1988}. Fritsch e Franco aplicam o problema à experiência brasileira e destacam que movimentos de ouro exerceram forte efeito sobre crédito e atividade \cite{fritschfranco1992}. Durante uma fase de exportações e capitais favoráveis, a compra do ouro pelo mecanismo de conversão ampliava a base monetária. Quando a conjuntura se invertia, a retirada de metal exigia contração. A mesma regra que acomodava a expansão importava o choque negativo para o sistema monetário doméstico.

O Brasil também não possuía, no período, um banco central com instrumentos semelhantes aos desenvolvidos nas principais praças financeiras. O Banco do Brasil intervinha no câmbio e no crédito, mas compartilhava o campo monetário com o Tesouro e, depois de 1906, com a Caixa de Conversão. A capacidade de esterilizar grandes entradas ou saídas de reservas era limitada. Por isso, credibilidade e disciplina não devem ser tratadas como benefícios sem custo. Elas organizavam expectativas e contratos, mas podiam tornar mais estreito o espaço de resposta a uma crise.

Essa literatura oferece categorias para analisar o caso, mas não será projetada diretamente sobre os jornais. Termos como credibilidade, regra contingente e assimetria são instrumentos do pesquisador. Os capítulos empíricos perguntarão se, e por meio de que vocabulário, os contemporâneos discutiram confiança, honra contratual, defesa do lastro, elasticidade, circulação, crédito e custo da estabilidade.

\section{A crise monetária da primeira década republicana}

A criação da Caixa foi precedida por uma trajetória marcada pela expansão monetária do início da República, pela crise cambial e fiscal, pela guinada restritiva de 1898 e pela valorização da moeda nos primeiros anos do século XX. Essa sequência é indispensável porque os participantes de 1906 argumentavam à sombra das experiências recentes. Emissão, depreciação e instabilidade remetiam à década de 1890; equilíbrio orçamentário, resgate do papel e reputação externa remetiam aos governos de Campos Sales e Rodrigues Alves.

As reformas bancárias dos primeiros governos republicanos ampliaram direitos de emissão e facilitaram a expansão do crédito. Entre 1890 e 1891, a base monetária cresceu, a procura por importações se elevou e a taxa de câmbio sofreu forte depreciação. A interrupção de capitais, a redução das receitas de exportação e a instabilidade política agravaram o balanço de pagamentos. A desvalorização aumentava em mil-réis a receita dos exportadores, mas também encarecia importações e elevava o custo interno das despesas públicas realizadas no exterior \cite{fritsch1990,francolago2011}.

No café, a combinação de crédito abundante, imigração, ferrovias e câmbio depreciado favoreceu a expansão das plantações. Como novos cafezais exigiam tempo para entrar em produção, o aumento da oferta chegou ao mercado com atraso. A partir de meados da década, colheitas crescentes e preços internacionais em queda transformaram a crise do café em problema nacional. A compensação cambial preservava parte da receita em moeda doméstica, mas não eliminava a sobreprodução nem os efeitos fiscais e externos da depreciação \cite{saes1981,holloway1978}.

Tentativas de saneamento monetário foram adotadas e interrompidas ao longo da década. A guerra civil no Sul, as dificuldades bancárias e a necessidade de financiar o Tesouro limitaram a continuidade de uma contração. Em 1895 e 1896, déficits externos, queda do café e retração dos capitais anunciaram uma nova crise. O governo precisava pagar dívida e outras obrigações em moeda estrangeira, mas arrecadava em mil-réis e dependia fortemente dos direitos de importação. A queda do câmbio multiplicava o custo doméstico do serviço externo ao mesmo tempo em que enfraquecia a base fiscal.

\section{O \textit{funding loan} e a guinada ortodoxa}

O acordo negociado em 1898 com os Rothschild reorganizou temporariamente o serviço da dívida externa. O \textit{funding loan} permitiu que juros fossem pagos por meio de novos títulos e adiou amortizações, evitando desembolsos imediatos que o governo já não conseguia sustentar. Em contrapartida, vinculou o refinanciamento a medidas de ajuste fiscal e monetário. Depósitos em moeda nacional associados ao serviço refinanciado foram retirados de circulação e incinerados, contribuindo para a contração do papel-moeda \cite{fritsch1980,fritsch1990,francolago2011}.

O acordo não deve ser reduzido a uma ordem externa aceita passivamente, nem apresentado como decisão puramente doméstica. O governo brasileiro buscava restaurar sua capacidade financeira, reduzir o custo cambial de suas obrigações e recuperar acesso ao crédito. Os banqueiros, como agentes e credores, condicionavam o refinanciamento à proteção dos contratos. A relação era assimétrica, mas operava por negociação entre interesses que não eram idênticos. Esse ponto será retomado quando a imprensa nomear os Rothschild, os chamar de agentes financeiros ou discutir a autoridade exercida sobre recursos brasileiros.

Joaquim Murtinho, ministro da Fazenda de Campos Sales, tornou-se o principal representante da política de equilíbrio orçamentário, contração monetária e valorização cambial. O horizonte declarado era o retorno à paridade legal de 1846, de 27 pence por mil-réis. Como a taxa corrente se encontrava muito abaixo desse nível, a meta implicava uma valorização prolongada e queda dos preços internos. Fundos de resgate e de garantia deveriam retirar papel de circulação e acumular recursos para a futura conversibilidade \cite{neuhaus1975,fritsch1980}.

A política melhorou as contas públicas e contribuiu para restaurar a confiança financeira, mas seus efeitos internos foram severos. A contração de 1899 e 1900 pressionou bancos e atividade econômica. A crise bancária de 1900 levou o governo a intervir no Banco da República do Brasil, fornecendo recursos e reorganizando sua atuação. A experiência evidenciou a necessidade de algum mecanismo de liquidez e de maior capacidade pública para intervir no mercado de câmbio. A ortodoxia aplicada não eliminou a intervenção estatal; ao contrário, a estabilização ajudou a formar instrumentos que mais tarde seriam usados de maneira ativa.

Rodrigues Alves e Leopoldo de Bulhões preservaram o objetivo de longo prazo da conversibilidade ao par, mas abrandaram a intensidade da contração. A recuperação das exportações, o crescimento da borracha, os saldos comerciais e a volta dos capitais europeus favoreceram a apreciação do mil-réis. O Banco da República foi reorganizado como Banco do Brasil em 1905, sob controle governamental e com funções relevantes no mercado de câmbio \cite{fritsch1980,oliveirasilva2001}. Assim, quando a controvérsia de 1906 se abriu, a valorização já não era apenas um projeto doutrinário. Era um processo em curso, apoiado por uma conjuntura externa favorável e por instituições reconstruídas após a crise.

\section{Valorização cambial e custos distributivos}

Na cotação então utilizada, expressa em pence por mil-réis, uma elevação do número indicava valorização da moeda brasileira. O movimento reduzia, em mil-réis, o custo do serviço da dívida e das importações. Também podia aliviar o preço interno de bens importados e fortalecer a confiança dos credores. Para exportadores, porém, a mesma quantidade de moeda estrangeira passava a render menos moeda nacional. Quando os preços internacionais do café estavam deprimidos e a oferta continuava a crescer, a apreciação comprimia a renda doméstica do setor.

Os efeitos não podem ser divididos mecanicamente entre dois grupos homogêneos. Cafeicultores endividados, comissários, exportadores, bancos, importadores, consumidores urbanos, trabalhadores e governos estaduais ocupavam posições distintas. A defesa de uma taxa baixa podia expressar interesse cafeeiro, preocupação com emprego, preferência por maior circulação ou oposição ao ritmo da deflação. A defesa de taxa alta podia invocar contratos, custo de vida, finanças públicas ou o retorno ao padrão legal. Mesmo dentro de São Paulo, a organização da defesa do café envolveu conflitos sobre financiamento, repartição de riscos e papel da União \cite{perissinotto1999,torelli2004}.

Essa multiplicidade recomenda cautela com as categorias de ortodoxos e expansionistas. Elas são úteis para localizar polos, mas não devem fundir posição sobre o nível do câmbio, preferência monetária e pertencimento social. Um defensor da conversibilidade podia rejeitar a continuação da valorização até 27 pence. Um crítico da Caixa podia considerar a taxa de 15 baixa demais ou, ao contrário, duvidar de que a instituição fornecesse elasticidade suficiente. O mesmo ator podia defender estabilidade e divergir sobre os meios usados para alcançá-la.

\section{A crise cafeeira e o Convênio de Taubaté}

A deterioração dos preços do café estimulou propostas de intervenção desde o fim da década de 1890. São Paulo restringiu novos plantios em 1902, e a organização de um programa de valorização ganhou força à medida que se aproximava uma colheita excepcional. Em fevereiro de 1906, os presidentes de São Paulo, Minas Gerais e Rio de Janeiro celebraram o Convênio de Taubaté. O plano previa compras e retenção de café financiadas por empréstimo, medidas de controle da oferta e instrumentos fiscais. A estabilização cambial em nível conveniente ao setor exportador integrava a arquitetura original \cite{holloway1978,torelli2004,francolago2011}.

O convênio reunia políticas relacionadas, mas institucionalmente distintas. A compra de estoques procurava sustentar o preço externo do produto. A Caixa de Conversão deveria limitar a valorização do mil-réis e emitir notas contra ouro a uma taxa fixa. A resistência do governo Rodrigues Alves e de grupos contrários ao desenho valorizador levou à separação dos projetos. A política do café e a Caixa seguiram tramitações próprias, ainda que permanecessem ligadas pela conjuntura e pelos interesses que as haviam aproximado.

\section{Dois consensos incompletos e três taxas}

Oliveira e Silva sustentam que a estabilidade cambial era amplamente desejada e que a divergência recaía principalmente sobre o nível da taxa e os meios para preservá-la \cite{oliveirasilva2001}. Essa formulação é um ponto de partida forte para a pesquisa. Ela desloca a questão de uma escolha abstrata entre estabilidade e desordem para o conteúdo distributivo e institucional da estabilidade. Ao mesmo tempo, os jornais precisam mostrar até onde esse acordo se estendia, quem ficava fora dele e como seus termos mudaram entre 1906 e 1914.

Três números organizam a controvérsia inicial. A taxa de 27 pence por mil-réis era a paridade legal estabelecida no século XIX e continuava a orientar os defensores da valorização até o par. A taxa de 12 pence apareceu em propostas ligadas à defesa do café e nas emendas de Alcindo Guanabara. David Campista formulou o projeto de 15 pence, taxa incorporada pela lei de 6 de dezembro de 1906 \cite{torelli2004,brasil1906caixa}. A Caixa, portanto, não foi criada a 12 pence.

A diferença entre essas taxas condensava alternativas de política. Fixar a 12 significava interromper mais cedo a valorização e oferecer maior proteção em moeda nacional aos exportadores. Fixar a 15 reconhecia a taxa próxima ao mercado e produzia um compromisso entre estabilização, interesses do café e receios de expansão excessiva. Manter 27 como horizonte significava não cristalizar a depreciação acumulada desde o início da República. A escolha também afetava o volume de notas emitidas por unidade de ouro, os incentivos à entrada de metal e a distribuição dos custos da transição.

O desenho aprovado não tornou conversível todo o papel-moeda existente. A Caixa recebia ouro e certas moedas estrangeiras, emitindo notas próprias conversíveis à taxa fixada. O restante da circulação permanecia formalmente inconversível. Oliveira e Silva descrevem o resultado como um padrão-ouro na margem: a instituição estabelecia um limite para a valorização e o Banco do Brasil podia intervir para manter a taxa de mercado próxima da paridade, mas a capacidade de defender o limite inferior dependia das reservas e da conjuntura \cite{oliveirasilva2001}.

\interpretacaoprovisoria{A hipótese que orienta os capítulos seguintes é que, em 1906, a recusa do padrão-ouro já não organizava o principal campo de conflito. A disputa se concentrava no nível da taxa, no ritmo da valorização, nos limites da emissão conversível, nos fundos de resgate e garantia, nas competências institucionais e nos custos distributivos da estabilidade. Essa formulação será testada, e não presumida, na leitura comparada dos jornais.}

O capítulo seguinte examina como esses objetos apareceram durante o Convênio de Taubaté, a separação das propostas, a tramitação do projeto Campista e a aprovação da Caixa. A questão não será apenas qual jornal apoiou ou rejeitou a instituição, mas que Caixa cada voz aceitava, temia ou pretendia construir.

